\documentclass[
    hungarian,
    parspace
]{elteikthesis}

\usepackage[utf8]{inputenc}
\usepackage{minted}

% --- 1. NYELVI BEÁLLÍTÁS (Ez oldja meg a \codelabel hibákat!) ---
\documentlang{hungarian}

% --- 2. KÖTELEZŐ METAADATOK ---
% A sortörés (\\) megoldja a túl hosszú cím (Overfull \vbox) problémáját
\title{Transzpilált nyelv koncepcióigazolása \\ fordítási idejű invariánsokkal} 
\date{2026} % A védés éve (a sablon ezt várja a \thesisyear helyett)

\author{Hallgató Neve}
\degree{Programtervező informatikus BSc}

\supervisor{Témavezető Neve}
\affiliation{Egyetemi docens} % Témavezető beosztása

% --- 3. AZ EGYETEM ÉS KAR ADATAI (Ezek oldják meg a \facname hibát!) ---
\university{Eötvös Loránd Tudományegyetem}
\faculty{Informatikai Kar}
\department{Programozási Nyelvek és Fordítóprogramok Tanszék}
\city{Budapest}

% Ne felejtsd el az images/elte_cimer_szines.pdf fájlt a mappádban hagyni!
\logo{images/elte_cimer_szines} 
% --------------------------------------

\begin{document}

\maketitle
\tableofcontents


% ==========================================
% 1. BEVEZETÉS ÉS SZAKIRODALOM (Cél: 4-5 oldal)
% ==========================================
\chapter{Bevezetés}
Manapság az informatikának egyre több szektorában jelenik meg az az igény, hogy a kód ne csak hatékony, hanem bizonyíthatóan helyes is legyen. A beágyazott rendszerek, az orvosi szoftverek, vagy a pénzügyi tranzakciókat végző modulok esetében egyetlen szoftverhiba is katasztrofális következményekkel járhat. 

\section{A problémafelvetés}
Gyakran egyensúlyt kell keresnünk a program hatékonysága és a helyesség ellenőrzése között. Akárhányszor egy függvénynek elő- vagy utófeltételt szeretnénk megfogalmazni és azokat kódban ellenőrizni (például \textit{assert} utasításokkal), valamennyi futásidejű teljesítményt fel kell áldoznunk. Számos szituáció előfordul, ahol egy erősebb heurisztika megfogalmazásával hatékonyabb kódot tudnánk írni, de a futásidejű ellenőrzések lassítják a szoftvert. 

Ennek a problémának megoldásaként merült fel egy olyan imperatív programozási nyelv létrehozása, mely ezeket az elő- és utófeltételeket, invariánsokat ha lehetséges, akkor fordítási időben (compile-time) ellenőrzi. 

\section{Elméleti háttér}
% TODO: Írj 1-2 oldalt az alábbi témákról a terjedelem növeléséhez!
Ahhoz, hogy a fordítási idejű invariáns-ellenőrzést megvalósítsuk, több informatikai és matematikai koncepciót is ötvöznünk kell:
\begin{itemize}
    \item \textbf{Hoare-logika:} A programok helyességének formális bizonyítására szolgáló rendszer, amely előfeltételekből (preconditions), parancsokból és utófeltételekből (postconditions) áll.
    \item \textbf{Szimbolikus végrehajtás (Symbolic Execution):} A program nem konkrét értékekkel (pl. $x = 5$), hanem szimbolikus változókkal fut, így a program minden lehetséges végrehajtási útvonala (path) feltérképezhetővé válik.
    \item \textbf{SMT Megoldók (Z3):} A Satisfiability Modulo Theories (SMT) megoldók képesek bonyolult logikai formulák kielégíthetőségét vizsgálni. Projektünk a Microsoft által fejlesztett Z3 SMT megoldót használja a matematikai és memóriabiztonsági állítások igazolására.
\end{itemize}

% ==========================================
% 2. FELHASZNÁLÓI DOKUMENTÁCIÓ (Cél: 5-7 oldal)
% ==========================================
\chapter{Felhasználói dokumentáció}
A szoftver célközönsége elsősorban azon fejlesztők és kutatók csoportja, akik magas megbízhatóságú szoftvereket terveznek, és szeretnék a kódjukban lévő szerződéseket (contracts) fordítási időben verifikálni[cite: 31, 34].

\section{Telepítés és Rendszerkövetelmények}
A fordítóprogram és az analizátor futtatásához a következő szoftveres környezet szükséges[cite: 35]:
\begin{itemize}
    \item \textbf{GHC (Glasgow Haskell Compiler):} A forráskód lefordításához (ajánlott verzió: 9.x).
    \item \textbf{Cabal vagy Stack:} A Haskell csomagfüggőségek (pl. Megaparsec, Z3) kezeléséhez.
    \item \textbf{Z3 SMT Solver:} A gazdagépen globálisan telepítve kell lennie a Z3 binárisnak és a C API-nak (pl. Ubuntu alatt \texttt{apt install z3 libz3-dev}).
\end{itemize}

Az indítás a \texttt{cabal run} paranccsal történik, mely paraméterként várja a bemeneti forrásfájlt[cite: 36].

\section{A nyelv szintaxisa és szemantikája}
A létrehozott nyelv egy erősen típusos, C-szerű imperatív nyelv. Támogatja az alapvető típusokat (\texttt{Int}, \texttt{Double}, \texttt{Bool}, \texttt{Char}), valamint a mutatókat (\texttt{Ptr}), tömböket (\texttt{Array}), és a komplex adatszerkezeteket (\texttt{struct}, \texttt{union}, \texttt{enum}).

\subsection{Függvények és Szerződések (Contracts)}
A nyelv legfőbb újítása a függvény-szignatúrákba ágyazott szerződések rendszere. Egy függvény deklarációja az argumentumok után opcionálisan tartalmazhat egy előfeltétel-blokkot és egy utófeltétel-blokkot:

\begin{minted}{c}
Int divide(Int a, Int b)(b > 0)(res != 0) { 
    return a / b; 
}
\end{minted}
A fenti kódban a \texttt{(b > 0)} a hívó félre vonatkozó előfeltétel, míg a \texttt{(res != 0)} a függvény által garantált utófeltétel (ahol a \texttt{res} kulcsszó a visszatérési értéket jelöli).

\subsection{Ciklus invariánsok}
A \texttt{while} ciklusok elemzése szimbolikus úton gyakran végtelen állapottérhez vezet. Ezt a nyelv a ciklusinvariánsok explicit megadásával kezeli:
\begin{minted}{c}
while (i < 10) (i >= 0) {
    i = i + 1;
}
\end{minted}

\section{Hibaüzenetek és diagnosztika}
Amikor a statikus analizátor egy szerződés megsértését vagy egy memóriahibát (pl. tömb alul/túlindexelés) észlel, a fordítás leáll, és a rendszer a következő formátumú diagnosztikai üzenetet adja[cite: 41]:
\texttt{fájlnév:sor:oszlop - Hiba típusa: Kifejezés}

Például, ha a fenti \texttt{divide} függvényt nullával hívjuk meg:
\begin{minted}{text}
test-input:5:5 - Precondition Failed in: divide: y > 0
\end{minted}

% ==========================================
% 3. RENDSZERTERV (Cél: 8-10 oldal)
% ==========================================
\chapter{Megoldási terv (Rendszerterv)}
Ez a fejezet bemutatja az alkalmazás architekturális felépítését, az adatfolyamot és a memóriamodellt[cite: 45].

\section{Architektúra és Adatfolyam}
% TODO: Ide illessz be egy UML Komponens diagramot! [cite: 48, 53]
A rendszer három fő logikai rétegre bontható[cite: 47]:
\begin{enumerate}
    \item \textbf{Lexer és Parser:} Beolvassa a nyers szöveget, ellenőrzi a szintaktikai szabályokat, és felépíti az Absztrakt Szintaxisfát (AST).
    \item \textbf{Analyzer (Hibrid Állapotgép):} Bejárja az AST-t. Ahol lehetséges, konkrét kiértékelést (Concrete Evaluation) végez, ahol nem (pl. ismeretlen szenzoradatok, függvényparaméterek), ott szimbolikus végrehajtásra tér át.
    \item \textbf{Z3 SMT Backend:} A szimbolikus formulákat lefordítja a Z3 SMT megoldó nyelvére, és matematikai bizonyításokat végez az invariánsokon és biztonsági szabályokon.
\end{enumerate}

\section{Az Absztrakt Szintaxisfa (AST) szerkezete}
A nyelv szintaxisát a \texttt{Parser.hs} modulban definiált ADT-k (Algebraic Data Types) írják le[cite: 50]. 

\begin{minted}{haskell}
data Expr = Expr SourcePos ExprNode
\end{minted}
Minden kifejezés (\texttt{Expr}) egy \texttt{SourcePos} csomagolóval van ellátva. Ez kritikus a pontos hibaüzenetek generálásához az Analizátor fázisban.

A memóriakezelést és a végrehajtás-vezérlést a \texttt{StmtNode} írja le, amely magában foglalja a deklarációkat, értékadásokat, és a vezérlési szerkezeteket (\texttt{If}, \texttt{While}, \texttt{For}).

\section{Memóriamodell és Állapotkezelés}
A szimbolikus elemző működéséhez elengedhetetlen egy robusztus memóriamodell[cite: 49]. Az állapotot az \texttt{AnalyzerState} rekord fogja össze:
\begin{itemize}
    \item \texttt{SymEnv}: A változóneveket memóriacímekre (Addr) képezi le.
    \item \texttt{Heap}: Konkrét futásidejű értékeket (pl. \texttt{VInt 5}) tárol az adott memóriacímeken.
    \item \texttt{SymStore}: A szimbolikus végrehajtás során keletkező matematikai kifejezéseket, "képleteket" tárolja.
    \item \texttt{facts}: A végrehajtási útvonal (Path Condition) során összegyűjtött logikai állítások listája.
\end{itemize}

% ==========================================
% 4. MEGVALÓSÍTÁS (Cél: 10-12 oldal)
% ==========================================
\chapter{Megvalósítás}
Ebben a fejezetben a funkcionális programozási paradigma sajátosságait, a monadikus transzformátorokat és az algoritmusok részleteit tekintjük át[cite: 56]. A Haskell nyelv választása lehetővé tette a kód erős strukturáltságát és a magas megbízhatóságot[cite: 61].

\section{A Szintaktikai Elemző (Parser) Implementációja}
A parser a \texttt{Megaparsec} könyvtárra épül. A kifejezések elemzésére a \texttt{Control.Monad.Combinators.Expr} modul \texttt{makeExprParser} függvényét alkalmaztam, amely lehetővé tette egy precedencia-táblázat (\texttt{operatorTable}) definiálását:
\begin{minted}{haskell}
operatorTable :: [[Operator Parser Expr]]
operatorTable =
    [ [ prefix "-" (UnOp "-"), prefix "!" (UnOp "!") 
      , prefix "*" Deref, prefix "&" AddrOf ]
    , [ binary "*" (BinOp "*"), binary "/" (BinOp "/") ]
    , [ binary "+" (BinOp "+"), binary "-" (BinOp "-") ]
    ...
\end{minted}
Ez garantálja a C-nyelvi precedenciaszabályok pontos betartását, elkerülve a terjedelmes és nehezen olvasható rekurzív ereszkedő (recursive descent) kód megírását.

\section{A Szimbolikus Végrehajtó Motor (\texttt{symExec})}
A rendszer "lelke" a \texttt{symExec} algoritmus, amely a szimbolikus kiértékelést végzi. Az elágazások (\texttt{If}) kezelése különösen komplex:
Amikor a szimbolikus végrehajtó egy elágazáshoz ér, ahol a feltétel értéke fordítási időben nem dönthető el (ismeretlen), a motor \textbf{kettéágazik}. Párhuzamosan kiértékeli a \textit{Then} és az \textit{Else} ágakat. 

A két ág végén a memóriastruktúrák (\texttt{SymStore}) eltérhetnek. Az egyesítést (Merge) egy Ternáris operátor generálásával oldottam meg, amely a feltétel függvényében választja ki a megfelelő értéket (SSA - Static Single Assignment formához hasonlóan).

\section{Biztonsági Ellenőrzések (Safety Extraction)}
Minden egyes kifejezés kiértékelése előtt lefut az \texttt{extractSafety} függvény, amely implicit előfeltételeket generál[cite: 67]. Például osztás operátor detektálása esetén automatikusan generál egy \texttt{b != 0} ellenőrzést, vagy tömb indexelés esetén a Z3-ba elküldi az alsó és felső indexhatárok ellenőrzését:
\begin{minted}{haskell}
(pos, "Array Index Out of Bounds", genExpr (BinOp "<" idx (genExpr (IntLit sz))))
\end{minted}

\section{Z3 Integráció és a \texttt{proveZ3} függvény}
Az SMT megoldóval való kommunikáció a \texttt{Z3.Monad} interfészen keresztül történik. A \texttt{translateZ3} függvény az AST-t Z3 sortokká (típusokká) és formulákká konvertálja. A bizonyítás trükkje: nem magát a feltételt akarjuk igazolni (mert az SMT a kielégíthetőséget -- SAT -- vizsgálja), hanem a feltétel \textbf{tagadását} (negation) adjuk a megoldónak. Ha a tagadás \textit{Unsat} (kielégíthetetlen), akkor a tétel bizonyítottan minden lehetséges esetben \textbf{Igaz}.

% ==========================================
% 5. TESZTELÉS (Cél: 5-8 oldal)
% ==========================================
\chapter{Tesztelés}
A szoftver robusztusságának és helyességének bizonyítása érdekében kiterjedt tesztelési keretrendszert építettem ki a \texttt{Hspec} modul használatával[cite: 70].

\section{Parser Modultesztek}
A parser tesztelése során feketedoboz tesztelést alkalmaztam. Az elemzett AST fát egy speciális \texttt{clearPos} függvény segítségével megtisztítottam a \texttt{SourcePos} adatoktól, így lehetővé vált a tiszta strukturális összehasonlítás.
Teszteltem a komplex mutatókonfigurációkat (pl. \texttt{const Int * const p;}), a precedenciaszabályokat, és a negatív teszteket is (olyan kódok, amiknek szintaktikai hibával el kell hasalniuk).

\section{Analizátor és Z3 Tesztesetek}
Az analizátor tesztelése a legfontosabb, hiszen itt bizonyosodik be a szoftver hozzáadott értéke. A tesztesetek az alábbi kategóriákat fedik le:
% TODO: Másold be ide a ParserTest.hs "Compile-Time Symbolic Analyzer" részéből a legérdekesebb teszteket (mint kódblokk) és magyarázd el a működésüket hosszan! [cite: 74, 75]

\subsection{Implicit utófeltételek felfedezése}
A motor képes a "sandbox" végrehajtás segítségével kitalálni a paraméterek módosulásait és a visszatérési értékek összefüggéseit anélkül, hogy a felhasználó ezt külön annotálta volna:
\begin{minted}{c}
Int add(Int a, Int b) { return a + b; }
Void process(Int val, Int x, Int y)(val == x + y) { return; }
\end{minted}
A fenti kódban a \texttt{process} függvény sikeresen lefut, mert a fordító automatikusan generált egy \texttt{res == a + b} utófeltételt az \texttt{add} függvényhez.

\subsection{Tömb határellenőrzések (Array Bounds Checking)}
A rendszer Z3 segítségével bizonyítani tudja az indexek helyességét. Ha az elemző a \texttt{arr[10] = 99;} sort látja egy 5 elemű tömbre, sikeresen dobja a \texttt{Array Index Out of Bounds: 10 < 5} fordítási idejű hibát.

\subsection{Ciklusok feloldása (Bounded Model Checking)}
A tesztek bizonyítják, hogy a kód képes automatikusan "kigöngyölni" a ciklusokat (loop unrolling). Egy 3 iterációs ciklust explicit invariánsok megadása nélkül is képes matematikailag bizonyítani a rendszer, amely jelentősen növeli a felhasználói élményt és csökkenti a fejlesztőre háruló teher mértékét.

% ==========================================
% 6. ÖSSZEGZÉS ÉS SZOFTVER FUTTATÁSA (Cél: 2-3 oldal)
% ==========================================
\chapter{Összegzés és Értékelés}
A bemutatott projekt keretében sikeresen elkészült egy olyan koncepcióigazoló (Proof of Concept) compiler technológia, amely egyesíti a hagyományos imperatív nyelvek szintaxisát a formális bizonyító rendszerek erejével.

\section{Szoftver futtatása és Robusztusság}
A tesztek és a kézi futtatások alapján a program a tervnek megfelelően működik[cite: 79]. Robusztussága kiemelkedő: hibás bemenet, végtelen ciklust eredményező forráskód, vagy ellentmondásos szimbolikus állítások esetén sem omlik össze ("száll el"). Az erőforrásokkal jól gazdálkodik (a \texttt{fuel} alapú limitált szimbolikus fa-bejárás garantálja, hogy a futásidő nem nő exponenciálisan a végtelenségig)[cite: 81, 83].

\section{További fejlesztési lehetőségek}
A jövőben a nyelv tovább bővíthető dinamikus memóriakezeléssel (pl. \texttt{malloc/free}), amelynek verifikálása a Separation Logic (szeparációs logika) beépítését igényelné. Szintén hasznos lenne egy kódgeneráló (Backend) modul írása, amely a megbízhatónak ítélt AST-t C vagy LLVM IR (Intermediate Representation) kódra fordítaná, így a lefordított program ténylegesen futtathatóvá válna operációs rendszer szinten is.

% TODO: Irodalomjegyzék hozzáadása BibTeX segítségével (Z3 papírok, Megaparsec doksi, Haskell könyvek)! 

\end{document}